\section{Complete Behavioral Tables and Model Exclusions}
\label{app:behavior}

\subsection{Requested response-surface panel on the full grid}

\begin{table}[H]
\centering
\small
\caption{Preliminary numerical exact accuracy over the full \(10\times5\times10\) grid (500 requests per available model--mode cell). These descriptive values precede the matched-policy rerun.}
\label{tab:full-grid-panel}
\begin{tabular}{@{}lrrrr@{}}
\toprule
Model / logical slot & Direct & Index & Bullet & CoT \\
\midrule
Qwen3-4B                  & .432 & .888 & .806 & .928 \\
Qwen3-8B                  & .570 & .962 & .894 & .956 \\
Qwen3-14B                 & \multicolumn{4}{c}{\textit{not yet evaluated}} \\
Qwen3-32B                 & .600 & .974 & .858 & .976 \\
Gemma4-E4B                & .248 & .882 & .796 & .832 \\
GLM-4 / GLM-Z1\({}^\dagger\) & .302 & .710 & .636 & .684 \\
Cogito-v1-Preview-8B      & .376 & .516 & .776 & .720 \\
Nemotron-Nano-v2-9B       & .236 & .666 & .580 & .636 \\
Nemotron-3-Nano-4B        & .366 & .652 & .546 & .654 \\
\bottomrule
\end{tabular}
\vspace{0.25em}

\parbox{.94\textwidth}{\footnotesize \({}^\dagger\)The GLM row stitches three modes from GLM-4-9B and CoT from GLM-Z1-9B; it is excluded from same-checkpoint mode effects. Gemma4-E4B Direct numerical exactness (.248) differs from registered success (.240); native CoT numerical exactness (.832) differs from registered success (.792).}
\end{table}

\subsection{Models excluded from the main panel}

Models are not removed because they disagree with the hypothesis. They are excluded from the main response-surface or mechanism panel when a frozen evaluation gate fails: incomplete mode coverage, a different checkpoint for the reasoning condition, severe parse/format/truncation confounding, or accuracy too close to floor for the planned causal contrast. Their numerical accuracy remains visible.

\begin{table}[H]
\centering
\scriptsize
\caption{Other audited models and why they are assigned to the appendix. Entries are full-grid numerical exact accuracy; a dash indicates that the mode was not run on that checkpoint.}
\label{tab:excluded-models}
\begin{tabularx}{\textwidth}{@{}P{.22\textwidth}rrrrY@{}}
\toprule
Model / checkpoint setup & Direct & Index & Bullet & CoT & Main-panel exclusion reason \\
\midrule
Qwen3-1.7B
  & .196 & .288 & .396 & .602
  & Low-capability floor; enumeration also has substantial parse/truncation failures. Retained as a phase-boundary sensitivity. \\
Gemma4-12B
  & .514 & .950 & .746 & .908
  & Native numerical accuracy is high but registered success is .314 because of final-channel format behavior. Not selected for the co-primary 8B-scale mechanism comparison. \\
DeepSeek-R1-0528-Qwen3-8B
  & -- & -- & -- & .672
  & Native-only cell; no within-checkpoint mode quartet. \\
Granite-3.3-Instruct-8B
  & .394 & .440 & .374 & .440
  & Low and weakly separated modes plus format/protocol failures; unsuitable for non-floor mechanism mediation. \\
Ministral Instruct / Reasoning
  & .658 & .892 & .756 & .770
  & CoT uses a separate reasoning checkpoint; native registered success is .020 despite .770 numerical exactness, creating a severe format confound. \\
OLMo3 Instruct / Think
  & .234 & .238 & .302 & .508
  & CoT uses a separate checkpoint; Native parse is .746 and truncation .254, while Direct is near floor. \\
\bottomrule
\end{tabularx}
\end{table}

The two Nemotron models remain in the requested cross-family panel, but their full per-model Index/Bullet surfaces are appendix material because enumeration has nontrivial parse/truncation failures. Nemotron-3 Native is comparatively clean. Cogito Index likewise requires parser/protocol repair before conditional signed-bias fits can be interpreted.

\section{Registered Experiment Ledger}
\label{app:ledger}

\begin{table}[H]
\centering
\scriptsize
\caption{Unified TODO ledger. The required output is part of the
preregistration, not a suggested visualization.}
\begin{tabularx}{\textwidth}{@{}P{.07\textwidth}P{.19\textwidth}P{.36\textwidth}Y@{}}
\toprule
ID & Experiment & Required output and controls & Claim enabled \\
\midrule
B0 & Configuration and failure audit
  & Frozen revisions/templates/budgets/parsers; request-level provenance,
  parse, truncation, protocol, and leakage flags
  & Comparable behavioral panel \\
B1 & Paired behavior grid
  & Four modes on 500 frozen stimuli; full requested model panel; exact count,
  registered success, error budget, and family-cluster intervals
  & Robust behavioral phenomenon \\
B2 & Held-out response surfaces
  & Qwen within-family development; untouched 14B prediction; grouped
  held-out scoring and coefficient stability
  & Optional empirical law \\
D1 & Controlled mechanism corpus
  & Tokenizer-specific equal-length siblings; stratified slots;
  family-disjoint splits; audited spans and invariant delimiters
  & Identifiable representation and patching tests \\
D2 & Prompt-policy controls
  & \(2\times2\) instruction by reasoning switch; predicate/policy placement;
  prefix hashes; paired decoding seeds; null scratchpad
  & Controlled algorithmic contrast \\
R1 & Registered-state decoding
  & Family-weighted held-out probes; nuisance controls; individual-state PCA;
  conditional dispersion and error strata
  & Count/progress representation \\
R2 & Direct dynamics and patches
  & Prefix-shift \(0/1/0\); total shift; same-position unchanged donors;
  predicate-after negative control; terminal tests
  & \(H_{\rm run}\) versus \(H_{\rm agg}\) \\
R3 & CoT anti-leakage and state control
  & No-index/random/permuted labels; new/non-target/repeated prefixes;
  residual versus trace/KV transport; steering and natural patches
  & Semantic progress carrier \\
H1 & Frozen role screen
  & \(B,C^B,T,D,U\); discovery/validation/test split; family-wise
  max-statistic; raw and conditional mass
  & Candidate functional components \\
H2 & Functional component tests
  & Query-local ablation; single-donor clean patch; path patch; semantic
  proximal endpoints and matched controls
  & Necessity and partial sufficiency \\
M1 & Integrated Qwen causal chain
  & Route/state/terminal factorial; coverage, progress, terminal, and complete
  answer endpoints
  & Mechanism-to-behavior link \\
M2 & Controls and Gemma replication
  & Common evidence; null scratchpad; semantic corruptions; independently
  discovered co-primary replication
  & CoT explanation and scope \\
S1 & Synthetic emergence
  & Frozen v20 basis; at least three seeds; all-sequence and objective-switch
  controls; fixed bundle tracked across checkpoints
  & Developmental support \\
R0 & Reproducible analysis registry
  & Versioned YAML plus manifests/checksums
  & All figures and tables reconstruct from frozen inputs and declared exclusions \\
\bottomrule
\end{tabularx}
\end{table}

For the Qwen mechanism claim, B0--B1, D1--D2, R1--R3, H1--H2, and M1 are
required. M2 is additionally required for the stated Qwen--Gemma functional
replication claim. B2 and S1 determine optional law and developmental scope;
R0 is required for every released result.

\section{Prompt Serialization and Activation Manifest}
\label{app:manifest}

The frozen V2 user message is serialized in this order:
\begin{verbatim}
You will need to count all city-score audit records
in the passage below.
A city-score audit record names one city and gives
that city's numeric score.
<passage>
[byte-identical frozen passage]
</passage>
How many city-score audit records are in the passage?
[mode-specific response-policy block]
\end{verbatim}
The response-policy block is Direct, Bullet, Index, or Native CoT as specified
in \cref{tab:prompt-modes}. Mechanism contrasts store both raw message content
and the output of the exact model-specific chat template; rendered text and
token IDs, rather than the unrendered message alone, determine prefix
identity.

Each activation row stores
\texttt{family\_id}, model and tokenizer revisions, interface,
predicate/policy placement, raw messages, rendered-prefix hash, complete input
and continuation token IDs, source character spans, model-specific source
token spans, registered landmark indices, data seed, decoding seed,
\texttt{source\_version}, hook name, and capture dtype. Assertions verify
byte-identical passages across interfaces, identical rendered token prefixes
whenever a shared-prefix claim is made, equal sequence length for each
clean/corrupt pair, and no truncation before a registered endpoint.

\section{Predicate and Policy Placement Matrix}
\label{app:placement}

\begin{table}[H]
\centering
\small
\caption{The causal mask and prompt serialization determine which state
comparisons are identifiable. Only the first row belongs to the frozen
behavioral grid; the remaining rows are mechanism controls.}
\begin{tabularx}{\textwidth}{@{}P{.23\textwidth}P{.29\textwidth}Y@{}}
\toprule
Condition & Passage-time information & Valid use \\
\midrule
Predicate-before, suffix-after
  & Predicate known; response suffix unseen; reasoning policy also unseen
  only when the rendered-prefix audit passes
  & Shared \(H_i\) candidate when prefix identity passes; otherwise analyze
  within interface; mode-specific \(Z_D,P_k\) remain valid \\
Predicate-after
  & Task predicate unseen
  & Content/position control at \(H_i\); terminal tests remain valid \\
Predicate-repeat
  & Predicate known and refreshed after passage
  & Separates passage-time availability from instruction recency \\
Policy-before
  & Predicate and response policy known
  & Tests interface-specific passage encoding; mechanism only \\
\bottomrule
\end{tabularx}
\end{table}

\section{Planned Result Figures and Reproducibility Checklist}

The planned empirical figures are:
\begin{enumerate}
  \item behavior surfaces and the complete failure budget (TODO B1--B2);
  \item registered-state geometry, dynamic updates, and natural-state patches
  (TODO D1--D2, R1--R3);
  \item real attention maps beside held-out ablation and patch effects for the
  three component roles (TODO H1--H2);
  \item intervention-to-coverage-to-progress-to-answer effects, including the
  common-evidence control (TODO M1--M2); and
  \item synthetic emergence order with seed uncertainty (TODO S1).
\end{enumerate}

The final artifact freezes model and tokenizer files, chat templates,
generation parameters, separate data and decoding seeds, stimulus and
rendered-prefix hashes, parser and alignment code, character-to-token maps,
landmark definitions, hook locations, intervention code, checkpoint IDs, and
all discovery/validation/test splits. Request-level outputs are retained
rather than only aggregates. Hardware, precision, batching, library versions,
and synthetic generator provenance are reported.

\section{Prompt, Trace, and Intervention Controls}
\label{app:controls}

\begin{itemize}
  \item Direct total-only; numbered list; bullets; index-free item tags; native CoT; equal-length semantically null scratchpad.
  \item Same faithful list in user versus assistant prefix; model-generated list replayed under teacher forcing; Direct with a prefilled ledger.
  \item Correct identity with wrong attribute; wrong identity with correct
  attribute; missing, invalid, and same-template distractor records. A second
  independent valid input occurrence still contributes \(+1\); only repeated
  retrieval of the same occurrence ID is a CoT progress no-op.
  \item Candidate ordinal fixed while validity flips; number of seen candidates fixed while valid-prefix count changes.
  \item Source order permutation, visible-index permutation, symbolic indices, and item anchors before versus after visible numerals.
  \item Fixed trace extent with different unique coverage; fixed coverage with different duplicate counts; fixed answer position with variable trace extent.
  \item Fixed \((N,L)\) with dispersed versus clustered targets and early/middle/late source-depth strata.
  \item Same final total with different source sets and traces; same trace length with different totals.
  \item Query-local mean/resample ablation and zero-ablation sensitivity;
  primary rescue with one matched natural donor; score-matched, random, and
  wrong-layer bundles; wrong query rows and source donors; norm-matched
  subspaces.
  \item Complete parsed numeric-string probability and free-running exact count, never only a first-digit logit.
\end{itemize}

\section{Additional Statistical Specification for State Dispersion}
\label{app:noise}

For interface \(m\), registered site \(s\), layer \(\ell\), and semantic label
\(c\), fit a family-weighted decoder on discovery families only. On locked
families define decoder bias, variance, and mean squared error:
\begin{equation}
b_{ms\ell}(c)=\E[\widehat c-c\mid c],
\qquad
v_{ms\ell}(c)=\Var(\widehat c-c\mid c),
\qquad
q_{ms\ell}(c)=v_{ms\ell}(c)+b_{ms\ell}(c)^2.
\label{eq:decoder-bias-var}
\end{equation}
Decoder error and \(\nu_{\rm state}\) in \cref{eq:state-snr} answer different
questions and are reported separately. The former measures accessibility to a
chosen decoder; the latter measures within-label dispersion relative to
adjacent separation after nuisance residualization.

On the controlled \(N\in\{0,\ldots,6,8,10\}\) corpus, estimate prespecified
slopes of \(\log(\nu_{\rm state}+\epsilon)\),
\(\log(v+\epsilon)\), and \(\log(q+\epsilon)\) against count or semantic
progress while controlling for context length, depth, terminal position,
trace extent, template, and correctness stratum. Nonadjacent sparse values use
their observed \(\Delta c\) and are never treated as missing adjacent counts.
Family-cluster intervals resample families. The fair interface comparison is
Direct \(Z_D\) versus controlled CoT \(Z_C^{\rm ctrl}\); \(H_i\) versus \(P_k\)
is cross-role and exploratory. A larger-count-noise claim requires convergent
separation, dispersion, decoder-error, and causal-transport evidence rather
than one significant slope.

\section{Architecture-Specific Intervention Notes}
\label{app:architecture}

Both co-primary models use grouped-query attention, so an intervention must specify whether it acts on a query head, a shared key/value head, or the concatenated pre-output-projection head slice. For each layer, the artifact records tensor shape, rotary position convention, attention backend, and hook location. The causal decomposition is:
\[
\begin{aligned}
\text{attention pattern}
&\rightarrow
\text{source-weighted value}
\rightarrow
\text{query-head output before }o_{\rm proj}\\
&\rightarrow
\text{attention block output}
\rightarrow
\text{post-layer residual}.
\end{aligned}
\]
Pattern-only recovery supports routing; value/head-output recovery supports information transport; residual-only recovery localizes the effect no more finely than the layer. MLP output is patched separately only after a block-level state effect is established. Cross-mode and cross-model activation patches are exploratory because the residual bases and manifolds need not align.

\section{Claim-to-Evidence Reporting Rules}
\label{app:rules}

\begin{enumerate}
  \item Every exact number in the main text must trace to a frozen CSV/table and a manifest, not only to HTML prose.
  \item PCA and centroids describe geometry. ``Counter'' additionally requires
  held-out decodability, the registered \(+1/0\) dynamic test, and downstream
  transport.
  \item ``Broad aggregation,'' ``targeted retrieval,'' and
  ``progress transition'' remain candidate labels until independent
  query-local ablation and single-donor clean restoration establish the
  proximal semantic endpoint.
  \item Teacher-forced effects establish local competence. Natural execution requires the same signed effect on free-running prefixes before the first error.
  \item Synthetic results establish measurement feasibility, not the mechanism of a pretrained LLM.
  \item A response surface becomes a law only after untouched-model prediction. Failed validation remains in the paper as a negative result.
\end{enumerate}
